
\section{Method}
\label{sec:method}

\subsection{Problem Formulation}
\label{sec:problem}

We study \emph{unsupervised} post-training of unified multimodal models, where the only training data are raw images. Let a pretrained unified model $\mathcal{M}_\theta$ expose two native interfaces: (i) visual understanding $\mathcal{U}: (\mathcal{I}, q)\mapsto a$ (answer a question $q$ about image $\mathcal{I}$), and (ii) image generation $\mathcal{G}: t\mapsto \hat{\mathcal{I}}$ (generate an image from text prompt $t$). The training pool is an unlabeled image set $\mathcal{D}=\{\mathcal{I}_1,\ldots,\mathcal{I}_N\}$ with no captions, QA pairs, or preference labels.

To preserve base capabilities and keep the procedure architecture-agnostic, we freeze the backbone parameters $\theta$ and train only lightweight role adapters:
\begin{equation}
    \phi=\{\phi_p,\phi_s,\phi_g\}, \qquad \theta'=\theta\oplus\phi,
    \label{eq:adapted_params}
\end{equation}
where $\phi_p$ parameterizes a \textbf{Proposer} $\pi_{\phi_p}(q\mid\mathcal{I})$, $\phi_s$ a \textbf{Solver} $\pi_{\phi_s}(a\mid\mathcal{I},q)$, and $\phi_g$ a \textbf{Generator} $\pi_{\phi_g}(\hat{\mathcal{I}}\mid t)$. For diffusion-based backbones such as BLIP3o, $\phi_g$ includes LoRA modules on both the language-model text-to-conditioning pathway and the frozen diffusion denoiser/DiT projections; the denoiser weights themselves remain frozen. Rewards use only model-native operations (QA, captioning, and text-conditioned generation) plus similarity computations from frozen representations, with no external judge.

Our goal is to jointly improve visual understanding and image generation using \emph{self-derived} rewards:
\begin{equation}
    \begin{aligned}
        \max_{\phi_p,\phi_s,\phi_g}\ \mathcal{J}(\phi)=
        &\ \lambda_u\,\mathbb{E}_{\tau_u\sim\pi_{\phi_p,\phi_s}}\!\left[R_p(\tau_u)+R_s(\tau_u)\right] \\
        &+ \lambda_g\,\mathbb{E}_{\tau_g\sim\pi_{\phi_s,\phi_g}}\!\left[R_g(\tau_g)\right],
    \end{aligned}
    \label{eq:joint_objective}
\end{equation}
subject to: (C1) no curated supervision, (C2) no external reward model, and (C3) controlled policy drift via KL regularization to frozen reference policies. Eq.~\eqref{eq:joint_objective} naturally separates into two loops: a \emph{visual understanding loop} (Proposer/Solver) and an \emph{image generation loop} (Generator evaluated by the Solver). Below we define the rewards $R_p$, $R_s$, and $R_g$.

\subsection{Framework Overview}
\label{sec:overview}

We instantiate all three roles with LoRA adapters~\cite{hu2022lora} on the same frozen backbone (Figure~\ref{fig:framework}). Training alternates between \emph{visual understanding steps}, which update the Proposer and Solver from self-consistency-based rewards, and \emph{image generation steps}, which update the Generator using the Solver as an internal evaluator. The central coupling mechanism is the Solver’s dual role: as the Solver becomes a stronger visual reasoner, its internal evaluations provide progressively sharper training signals for image generation.

\begin{figure*}[t]
  \centering
  \includegraphics[width=\linewidth]{figures/main-figure.pdf}
  \caption{Overview of our \emph{Proposer}--\emph{Solver}--\emph{Generator} self-evolving framework. We attach three lightweight LoRA adapters to a frozen backbone and train using only unlabeled images. Left (visual understanding): the \emph{Proposer} generates questions and the \emph{Solver} answers under prompt perturbations; self-consistency and Solver Token Entropy (STE) provide rewards targeting the \emph{Solver}'s competence frontier. Right (image generation): the \emph{Generator} synthesizes images from text prompts and the \emph{Solver} evaluates outputs via QA fidelity and cycle-consistent captioning, coupling both loops for progressive co-improvement without external labels or reward models.}
  \label{fig:framework}
\end{figure*}


\subsection{Self-Consistent Visual Understanding}
\label{sec:understanding}

We learn to ask and answer informative questions from unlabeled images using two complementary signals: (i) \emph{framing-level robustness}, which measures whether answers remain stable under equivalent instructions, and (ii) \emph{token-level uncertainty}, which remains informative when output-level agreement becomes degenerate.

For an image $\mathcal{I}$, the Proposer samples a question $q \sim \pi_{\phi_p}(\cdot \mid \mathcal{I})$. We query the Solver under $N$ fixed prompt framings $\{\rho_1,\ldots,\rho_N\}$ (templates that preserve semantics but vary wording), i.e., $a_i=\text{Solver}_{\phi_s}(\mathcal{I}, \rho_i(q))$ for $i=1,\ldots,N$.
Let $p_c$ be the empirical distribution over distinct answers, and $a^*$ the modal answer. We measure sample-level self-consistency by entropy:
\begin{equation}
    H_{\text{sc}} = -\sum_{c} p_c \ln p_c, \quad p_c = \frac{|\{i : a_i = c\}|}{N}.
    \label{eq:sample_entropy}
\end{equation}
Low $H_{\text{sc}}$ indicates robustness to rephrasing; high $H_{\text{sc}}$ indicates unstable reasoning.

\paragraph{Solver Token Entropy (STE).}
$H_{\text{sc}}$ can collapse (all framings yield the same answer) even when the model is uncertain, producing weak curriculum signals. To recover a continuous notion of difficulty, we compute a token-level uncertainty score from the Solver’s next-token distributions. Let $\mathbf{p}_t$ be the next-token distribution at decoding step $t$ of an answer with length $T$, and $H_t = -\sum_{v} p_{t,v} \ln p_{t,v}$ its entropy. We define STE as:
\begin{equation}
    \hat{H} = \max_{t \in \{1, \ldots, T\}} H_t,
    \label{eq:ste}
\end{equation}
and convert $\hat{H}$ into a stable difficulty score $d_{\text{ste}}\in[0,1]$ by taking its percentile rank within a rolling window of recent $\hat{H}$ values. This quantile normalization makes STE comparable across training, even as the Solver’s overall confidence distribution shifts.

\paragraph{Proposer and Solver optimization.}
The Proposer is trained to generate questions near the Solver’s competence frontier: questions that are neither trivial nor unsolvable. To this end, we compute a bounded self-consistency score $r_{\text{sc}}\in[-1,1]$ that combines (i) an \emph{adaptive target} favoring medium self-consistency entropy (a Gaussian centered at a running EMA target $\mu$) and (ii) a \emph{local informativeness} term that down-weights degenerate cases such as unanimous votes with large margins. To prevent reward degeneracy when $H_{\text{sc}}\!\approx\!0$, we add STE difficulty:
\begin{equation}
    R_p = w_{\text{ste}} \cdot d_{\text{ste}} + w_{\text{sc}} \cdot r_{\text{sc}},
    \label{eq:proposer_reward}
\end{equation}
which yields an automatic curriculum that tracks the Solver as it improves.

We define the Solver reward $R_s$ to encourage framing-robust answers: it rewards agreement with the modal answer $a^*$ across prompt perturbations, with small penalties for non-canonical outputs and excessive verbosity. In practice, we optimize both Proposer and Solver with KL-regularized policy gradients and EMA baselines (GRPO~\cite{shao2024deepseekmath} over groups of $K$ sampled questions for the Proposer; REINFORCE over answer tokens for the Solver), using an adaptive KL coefficient to control drift from the frozen reference model.

This design keeps the optimization stable while allowing the Proposer--Solver pair to co-evolve under purely internal signals.


\subsection{Image Generation Assessment via Internal Evaluation}
\label{sec:generation}

We improve generation by treating the Solver as an internal evaluator. Given a prompt $t$ for a source image $\mathcal{I}$ (obtained by having the Solver caption $\mathcal{I}$), the Generator produces a candidate image $\hat{\mathcal{I}}=\text{Generator}_{\phi_g}(t)$. We compute two complementary rewards.

\paragraph{QA fidelity scoring.}
We sample $M$ diagnostic questions about $\mathcal{I}$ (from the Proposer) and obtain reference answers $a_j^{\text{ref}}$. We then ask the same questions on the generated image:
\begin{equation}
    S_{\text{fid}} = \frac{1}{M} \sum_{j=1}^{M} \mathbb{1}[\hat{a}_j = a^{\text{ref}}_j].
    \label{eq:fidelity}
\end{equation}
$S_{\text{fid}}$ measures how well semantic attributes are preserved from $\mathcal{I}$ to $\hat{\mathcal{I}}$.

\paragraph{Cycle-consistent captioning.}
QA fidelity is local; we add a global check by captioning the generated image to obtain $\hat{t}$ and measuring cycle consistency:
\begin{equation}
    S_{\text{cyc}} = \tfrac{1}{2}\,\text{sim}(f(t), f(\hat{t})) + \tfrac{1}{2}\,\text{sim}_{\text{vl}}(\hat{\mathcal{I}}, t),
    \label{eq:cycle}
\end{equation}
where $f(\cdot)$ is a frozen text embedding and $\text{sim}_{\text{vl}}$ is a frozen vision--language similarity.

\paragraph{Generator optimization.}
The Generator reward combines these signals, plus small diversity/contradiction terms ($S_{\text{div}}$ encourages diversity across multiple candidates; $S_{\text{ctr}}$ penalizes prompt--caption contradictions):
\begin{equation}
    R_g = w_{\text{fid}} \cdot S_{\text{fid}} + w_{\text{cyc}} \cdot S_{\text{cyc}} + w_{\text{div}} \cdot S_{\text{div}} - w_{\text{ctr}} \cdot S_{\text{ctr}},
    \label{eq:gen_reward}
\end{equation}
The update rule depends on the generator parameterization. For autoregressive UUG backbones that expose discrete image-token log probabilities, we update $\phi_g$ with the same KL-regularized policy-gradient recipe used in the visual understanding loop. For diffusion or flow-matching backbones such as BLIP3o, the generated image is produced by a continuous denoiser rather than by discrete image tokens; applying token-level policy gradients would train only text/caption traces and not the image-producing module. In this case we use a diffusion-compatible Generator update: reward-weighted denoising regression on LoRA modules in the denoiser/DiT and, when available, the text-conditioning adapter. The base backbone and denoiser weights remain frozen; $R_g$ only weights the adapter loss.

\paragraph{Coupling design.}
The two loops are coupled only through the Solver in the main method: the Proposer is trained purely from visual understanding signals, while image generation rewards become stronger as the Solver becomes a more reliable evaluator. We evaluate stronger generation-phase Proposer/Solver coupling only as ablations, not as the default training recipe.



% \paragraph{Prompt perturbation details.}
% The N=7 prompt framings are generated using a template-based system
% that varies: (1) instruction style (direct question vs. instruction
% vs. analysis request), (2) verbosity expectation (concise vs. detailed),
% (3) role framing (neutral vs. expert vs. student), and (4) output
% format (free-form vs. constrained choice). Each template preserves the
% core question content while changing surface-level phrasing.
%
% Example perturbations for question "What color is the car?":
%   rho_1: "What color is the car?"
%   rho_2: "Please analyze this image and identify the color of the car."
%   rho_3: "As a visual perception expert, what would you say is the
%           color of the car in this image?"
%   rho_4: "Look at the car in this image. What color is it? Give a
%           one-word answer."
%   rho_5: "Describe the color of the car shown."
%   rho_6: "I need to know: what color is that car?"
%   rho_7: "Examining this image carefully, determine the car's color."
%
% The key insight is that semantically equivalent but syntactically
% different framings probe whether the model's answer is grounded in
% visual understanding (consistent) or surface-level pattern matching
% (inconsistent). Temperature-based sampling probes stochastic
% uncertainty but not framing robustness.

% \paragraph{Bucket-stratified EMA baselines.}
% Standard REINFORCE uses a single EMA baseline b for variance reduction.
% We partition the reward space into 3 difficulty buckets based on
% d_ste quantiles: easy (0-0.33), medium (0.33-0.67), hard (0.67-1.0).
% Each bucket maintains its own EMA baseline with coefficient alpha=0.95.
% This prevents the common failure mode where a single baseline, dominated
% by easy questions, assigns artificially high advantage to medium
% questions, causing the Proposer to get stuck at intermediate difficulty.
% Bucket-stratified baselines ensure consistent gradient pressure within
% each difficulty tier.

% \paragraph{Verbosity penalty.}
% The Solver reward includes a verbosity penalty:
%   R_verb = -lambda_verb * max(0, len(answer) - L_target) / L_target
% where L_target is set to the 75th percentile of answer lengths within
% the current batch, and lambda_verb = 0.1. This prevents the Solver
% from producing excessively long answers to hedge uncertainty (e.g.,
% listing multiple possible answers). Without this penalty, the Solver
% tends to produce longer answers over training, which dilutes the
% self-consistency signal by introducing more opportunities for
% partial agreement.

% \paragraph{Adaptive KL scheduling.}
% The KL coefficient beta_A is adjusted every 100 steps:
%   if mean_KL > target_KL * 1.5:  beta_A *= 1.5
%   elif mean_KL < target_KL * 0.5:  beta_A *= 0.67
% This keeps the KL divergence from the reference model within a
% controlled range, preventing both collapse to the reference (too high
% beta) and uncontrolled drift (too low beta). The target KL is set
% to 0.1 nats for the Solver and 0.05 nats for the Proposer, reflecting
% the design principle that the Solver should be allowed more deviation
% (to improve capability) while the Proposer should stay closer to
% natural question distributions.

% \paragraph{Diversity bonus computation.}
% S_div is computed as the average pairwise LPIPS distance among the
% k=3 candidate images per prompt:
%   S_div = (2 / k(k-1)) * sum_{i<j} LPIPS(I_i, I_j)
% LPIPS uses the model's own visual encoder features rather than a
% separate perceptual network. This encourages the Generator to produce
% visually distinct candidates, preventing mode collapse to a single
% "safe" composition that satisfies QA fidelity but lacks variety.
% The diversity bonus weight (w_div=0.10) is deliberately small to
% avoid rewarding diversity at the expense of fidelity.
